\labsection{1}{Capture, filter, and read a live conversation}{sec:lab01-wireshark}

\begin{labproject}{Turn a packet capture into evidence}
\textbf{Coverage.} Chapter~\ref{ch:wireshark}.\\
\textbf{Materials.} \texttt{Lab01/capture\_checklist.md}, \texttt{Lab01/display\_filters.txt}.\\
\textbf{Goal.} Produce a capture that answers a specific question, rather than a screenshot that shows traffic existed.

\begin{labmode}
Capture on a network you are authorised to observe --- your own machine, your own lab segment, or the eNSP simulation. Capturing on a shared campus or corporate network without permission is a disciplinary matter and in many jurisdictions an offence.

Everything below works on loopback or your own interface, so there is no reason to point Wireshark at anyone else's traffic.
\end{labmode}

\medskip\noindent\textbf{Part A --- Ask the question first.}\par
The single most common mistake is starting the capture before deciding what it is meant to prove. You then face 40{,}000 packets and no criterion for which matter.

Write the question down before you press start. Good questions look like:

\begin{itemize}
  \item Did the DNS lookup for \cmd{example.com} succeed, and how long did it take?
  \item Did the TCP handshake complete, or was it reset?
  \item Which host sent the first packet of this conversation?
\end{itemize}

\begin{keyidea}
A capture filter decides what is \emph{recorded}. A display filter decides what is \emph{shown}. They use different syntaxes, and confusing them is the most common source of an empty result window.

Capture filter (BPF): \cmd{tcp port 80}\\
Display filter (Wireshark): \cmd{tcp.port == 80}

Capture broadly and filter afterwards. A capture filter that was too narrow cannot be widened once the traffic has passed.
\end{keyidea}

\medskip\noindent\textbf{Part B --- A disciplined capture.}\par

\begin{netsteps}
  \item Choose the interface carrying the traffic. If unsure, watch the sparklines in the Wireshark start screen.
  \item Start the capture, then generate the traffic --- not the other way round.
  \item Perform exactly one transaction. One page load, one ping sequence, one login.
  \item Stop the capture immediately. Ten seconds of relevant traffic beats ten minutes of noise.
  \item Save the raw \cmd{.pcapng} before filtering. That file is your evidence.
\end{netsteps}

\medskip\noindent\textbf{Part C --- Filters that earn their keep.}\par

\begin{commandmap}
\begin{tabularx}{\linewidth}{@{}>{\ttfamily}p{0.38\linewidth}X@{}}
\toprule
\rowcolor{MLPaleBlue!92}
\textrm{\bfseries Display filter} & \bfseries Shows \\
\midrule
ip.addr == 192.168.1.10 & Everything to or from one host \\
tcp.port == 80 || tcp.port == 443 & Web traffic in both directions \\
tcp.flags.syn == 1 \&\& tcp.flags.ack == 0 & Connection attempts only \\
tcp.analysis.retransmission & Packets Wireshark believes were resent \\
dns & Name resolution \\
icmp & Ping and unreachable messages \\
http.request.method == "GET" & Page requests \\
\bottomrule
\end{tabularx}
\end{commandmap}

\begin{pitfall}
\cmd{ip.addr != 192.168.1.10} does not mean what it appears to. A packet has both a source and a destination address, so the expression is true whenever \emph{either} differs --- which is nearly always.

Write \cmd{!(ip.addr == 192.168.1.10)} instead. The same trap applies to every address and port field.
\end{pitfall}

\medskip\noindent\textbf{Part D --- Read one TCP handshake end to end.}\par
Filter to a single conversation (right-click a packet, then \emph{Follow $\rightarrow$ TCP Stream}) and identify the three-way handshake:

\begin{lstlisting}[style=dcnterminal]
No.  Source        Destination   Proto  Info
1    192.168.1.10  93.184.216.34 TCP    49152 > 80 [SYN] Seq=0 Win=64240
2    93.184.216.34 192.168.1.10  TCP    80 > 49152 [SYN, ACK] Seq=0 Ack=1
3    192.168.1.10  93.184.216.34 TCP    49152 > 80 [ACK] Seq=1 Ack=1
\end{lstlisting}

Answer from your own capture: which host initiated, what ephemeral port did it choose, and how long elapsed between packets 1 and 3? That last figure is the round-trip time, and it is the number a network engineer actually cares about.

\begin{troubleshoot}
\textbf{No packets at all.} Wrong interface, or a capture filter with a typo. Clear the capture filter and try again.

\textbf{Traffic appears but not yours.} You are probably on a switched network seeing only broadcast and your own frames --- which is normal and correct.

\textbf{HTTPS shows no readable content.} That is TLS doing its job. You can still see the handshake, the SNI hostname, and the timing. Analyse those instead.
\end{troubleshoot}

\begin{tryit}
Capture a \cmd{ping} to an address that does not exist on your subnet. No reply arrives --- so what \emph{does} the capture contain, and what does that tell you about how the host tried to find the destination? Look for ARP.
\end{tryit}
\end{labproject}

\begin{verifybox}
\begin{itemize}
  \item The raw \cmd{.pcapng} is saved and submitted, not only screenshots.
  \item The question the capture answers is stated before the evidence.
  \item Display filters are quoted exactly as used.
  \item Packet numbers are cited when referring to specific packets.
  \item The capture was taken on a network you are authorised to observe.
\end{itemize}
\end{verifybox}

\begin{labdeliverable}
Submit the \cmd{.pcapng} file, the question it answers, the display filter used, an annotated screenshot of the three-way handshake with packet numbers, and the measured round-trip time. Add one sentence on why capturing broadly and filtering afterwards is the safer order.
\end{labdeliverable}
